\clearpage
\section{Calculs algébriques}
\begin{objectif}
Organiser un calcul long en blocs, expliquer les bornes d'une somme et choisir entre une simplification algébrique et un calcul numérique.
\end{objectif}
\subsection{Explorer : un atelier de mosaïque}
Une mosaïque comporte $n$ couronnes. La couronne de rang $j$, à partir de $j=0$, contient $4j+6$ carreaux. Dessiner les trois premières couronnes comme des rangées de points. Prédire le total pour huit couronnes, puis comparer une addition directe et un regroupement par paires.
Dans le diaporama, déplacer les bornes de la somme : quel terme entre ou sort ? Une expérience numérique aide à trouver une formule ; elle ne la démontre pas.
\subsection{Écrire et transformer}
\begin{definition}[indices]
Pour des entiers $p\leq q$, l'écriture $\sum_{j=p}^{q}a_j$ désigne l'addition des termes dont les indices parcourent les entiers de $p$ à $q$, bornes comprises. Le symbole $\prod$ remplace l'addition par une multiplication. Avec un intervalle d'indices vide, on convient que ces valeurs sont respectivement $0$ et $1$.
\end{definition}
Le nom de l'indice est local au calcul. Pour décaler les indices, on transporte à la fois les deux bornes et le terme : poser $r=j+h$ donne $j=r-h$. Écrire les correspondances aux deux extrémités avant de transformer la somme évite les termes perdus.
\begin{propriete}[progression arithmétique]
Une liste de $n\geq1$ termes, de premier terme $a$ et de pas $d$, a pour somme $na+d\,n(n-1)/2$.
\end{propriete}
\begin{demonstration}
Écrire la liste dans l'ordre direct puis dans l'ordre inverse. Chaque colonne donne $2a+(n-1)d$. Les $n$ colonnes constituent deux fois la somme cherchée ; diviser par deux donne la formule.
\end{demonstration}
\begin{methode}[annulations]
Dans une somme de différences $v_{j+1}-v_j$, écrire les premiers et derniers termes avant d'annuler les termes intérieurs. Dans un produit de quotients, vérifier d'abord que tous les dénominateurs sont non nuls.
\end{methode}
\subsection{Choisir des objets}
\begin{definition}[coefficient binomial]
Pour $0\leq k\leq n$, $\binom nk$ compte les choix de $k$ objets parmi $n$ objets distincts, sans ordre et sans répétition. Sa valeur est $n!/(k!(n-k)!)$, avec $0!=1$.
\end{definition}
En ordonnant les $k$ objets choisis, chaque choix est compté $k!$ fois ; les listes ordonnées sont au nombre de $n(n-1)\cdots(n-k+1)$. Cela justifie la formule. Dans $(a+b)^n$, choisir les $k$ facteurs qui fournissent $b$ donne
\[ (a+b)^n=\sum_{k=0}^n\binom nk a^{n-k}b^k. \]
Cette preuve explique les coefficients plutôt que de les faire mémoriser.
\subsection{Ateliers autonomes}
\begin{enumerate}
\item \textbf{Budget de carreaux.} Retrouver le total $M_n=\sum_{j=0}^{n-1}(4j+6)$ de la mosaïque. Calculer $M_8$, puis déterminer combien de couronnes complètes on peut construire avec 260 carreaux.
\item \textbf{Deux modes d'annulation.} Pour $n\geq1$, étudier $U_n=\sum_{j=1}^n1/((j+2)(j+3))$ et $V_n=\prod_{j=1}^n(j+3)/(j+2)$. Trouver une relation entre $U_n$ et $V_n$ sans effectuer les opérations une par une.
\item \textbf{Inventaire pondéré.} Un appareil possède cinq emplacements, chacun recevant soit une pièce de type A (trois modèles), soit une pièce de type B (deux modèles). Compter les configurations par le nombre de pièces B, puis vérifier le total sans utiliser de somme.
\end{enumerate}
\ifprof
\subsection{Préparer et discuter}
Commencer par les dessins de couronnes, puis faire défendre deux méthodes de comptage. Demander à un groupe de suivre les indices et à un autre de contrôler les unités. Pour l'inventaire, distinguer le choix des positions et le choix des modèles. Faire expliquer pourquoi une vérification pour huit couronnes ne suffit pas pour tout $n$.
\subsection{Solutions des ateliers}
\begin{enumerate}
\item Le regroupement donne $M_n=n(6+(4n+2))/2=2n(n+2)$. Donc $M_8=160$. La fonction $M_n$ croît pour $n\geq0$ ; $M_{10}=240\leq260<M_{11}=286$ : on peut construire 10 couronnes complètes.
\item L'identité $1/((j+2)(j+3))=1/(j+2)-1/(j+3)$ donne $U_n=1/3-1/(n+3)$. Dans le produit, les facteurs intérieurs s'annulent et $V_n=(n+3)/3$. Ainsi $U_n=(1-1/V_n)/3$.
\item Avec k pièces B, choisir leurs positions puis leurs modèles donne $\binom5k3^{5-k}2^k$. Les six effectifs sont $243, 810, 1080, 720, 240, 32$ et leur total vaut 3125. L'autre comptage donne directement $5^5=3125$, car chaque emplacement possède cinq choix.
\end{enumerate}

\fi
